% ==============================================================================
% =================================== Aufgabe 5 =================================
% ==============================================================================

\chapter{Aufgabe}
\label{chap:5 Aufgabe}
\thispagestyle{fancy}
\section*{Entwerfen Sie eine Übersicht der Prozesse, die durch das \acs{ERP}-System unterstützt werden. Mit welchen fünf Kennzahlen können Sie Auswirkungen des \acs{ERP}-Systems auf die Prozesse messen?} \par Ein \acs{ERP}-System fungiert als integratives Rückgrat der betrieblichen Informationsverarbeitung und unterstützt eine Vielzahl von Funktionsbereichen und Prozessen. \cite{109}\cite{110} Die folgende Übersicht fasst die wesentlichen durch ein \acs{ERP}-System unterstützten Prozesse zusammen, basierend auf der horizontalen und vertikalen Integration im Unternehmen: \cite{53}\cite{111} \par \textbf{Übersicht der durch das \acs{ERP}-System unterstützten Prozesse}
\begin{itemize}
    \item \textbf{Materialwirtschaft und Logistik:} Dies umfasst die \uline{Materialbedarfsplanung (\textit{\acs{MRP}})}, den \uline{Einkaufszyklus} (\textit{von der Bestellanforderung über die Lieferantenauswahl bis zur Bestellung}), die \uline{Bestandsführung} (\textit{Erfassung von Warenbewegungen in Echtzeit}) sowie die \uline{Lagerverwaltung} und die Rechnungsprüfung. \cite{111}\cite{112}\cite{113}\cite{114}\cite{40}
    \item \textbf{Produktion:} Das System unterstützt die \uline{Produktionsplanung} und \uline{-steuerung (\textit{\acs{PPS}})}, die Erstellung von \uline{Stücklisten} und \uline{Arbeitsplänen} sowie das Management von Fertigungskapazitäten sowohl in der Serien- als auch in der Einzelfertigung. \cite{115}\cite{116}\cite{117}\cite{118}\cite{119}
    \item \textbf{Vertrieb:} Hier werden Prozesse der \uline{Akquisition}, die Bearbeitung von Kundenanfragen und Angeboten, die \uline{Auftragsabwicklung}, die \uline{Verfügbarkeitsprüfung (\textit{\acs{ATP}})}, der Versand sowie die Fakturierung abgebildet. \cite{120}\cite{121}\cite{122}\cite{123}
    \item \textbf{Rechnungswesen und Finanzen:} \acs{ERP}-Systeme integrieren die \uline{Finanzbuchhaltung} (\textit{Hauptbuch, Debitor- und Kreditor-Buchhaltung}), die \uline{Anlagenbuchhaltung}, das \uline{Liquiditätsmanagement} sowie die \uline{Kosten-} und \uline{Leistungsrechnung} inklusive Kalkulationen und Controlling. \cite{56}\cite{124}\cite{125}\cite{126}\cite{127}
    \item \textbf{Personalwesen (\textit{\acs{HR}}):} Unterstützt werden die \uline{Personaladministration}, die \uline{Lohn-} und \uline{Gehaltsabrechnung}, die \uline{Zeitwirtschaft} (\textit{Zeiterfassung}) sowie die Personalbedarfsplanung und -entwicklung. \cite{128}\cite{129}\cite{130}\cite{131}
    \item \textbf{Querschnittsfunktionen:} Dazu zählen das \uline{Qualitätsmanagement} (\textit{Prüfplanung und -abwicklung}) sowie die \uline{Instandhaltung} von Betriebsmitteln. \cite{132}\cite{133}\cite{134}
\end{itemize}
\textbf{Fünf Kennzahlen zur Messung der Auswirkungen auf die Prozesse} \par Um die Effektivität und den Nutzen der \acs{ERP}-Einführung zu bewerten, können folgende Kennzahlen herangezogen werden:
\begin{enumerate}[label=\arabic*.), leftmargin=*]
    \item \textbf{Durchlaufzeit (\textit{\acs{z. B.} Bestellzyklus oder Auftragsabwicklungszeit}):} Durch die Automatisierung von Aktivitäten und die Eliminierung manueller Eingriffe können \acs{ERP}-Systeme (\textit{insbesondere im E-Procurement}) die \uline{Bestellzyklen um bis zu 70 \% verkürzen}. \cite{135}\cite{136}
    \item \textbf{Prozesskosten (\textit{Administrative Kosten}):} \acs{ERP}-Systeme zielen auf eine \uline{Reduzierung der administrativen Kosten} ab (\textit{oft zwischen 50\% und 70\%}), indem sie fragmentierte Abläufe standardisieren und den Aufwand für manuelle Datenpflege verringern. \cite{31}\cite{136}\cite{137}
    \item \textbf{Lagerumschlagshäufigkeit / Bestandskosten:} Eine verbesserte Disposition und Transparenz führt zu \uline{niedrigeren Lagerbeständen} und einer geringeren Kapitalbindung. Die Bestandskosten können durch optimierte Prognosemodelle signifikant gesenkt werden. \cite{135}\cite{136}\cite{137}\cite{138}
    \item \textbf{Termintreue (\textit{Liefertreue}):} \acs{ERP}-Systeme verbessern die \uline{Qualität logistischer Prozesse}, indem sie eine rechtzeitige Prüfloseröffnung und eine präzise Terminierung ermöglichen, was die Verlässlichkeit gegenüber dem Kunden steigert. \cite{136}\cite{139}\cite{140}
    \item \textbf{Datenqualität (\textit{Fehlerquote bei der Erfassung}):} Die Auswirkungen des Systems lassen sich an der \uline{Vermeidung von Erfassungs-} und \uline{Übertragungsfehlern} messen, da Daten nur einmal zentral erfasst werden und systemübergreifend konsistent zur Verfügung stehen. \cite{107}\cite{136}\cite{141}
\end{enumerate}

% ==============================================================================
% =================================== Aufgabe 5 =================================
% ==============================================================================